\section{Related Work} \label{sec:relatedWork}

\subsection{Requirements Ambiguity and Prompt Robustness}

Ambiguity and incompleteness have long been recognized as important causes of misunderstanding in software requirements. Prior studies in requirements engineering have investigated different forms of ambiguity, such as lexical, syntactic, semantic ambiguity, and proposed rule-based, statistical, or learning-based methods to detect ambiguous requirements~\cite{berry2003ambiguity,ezzini2022anaphoric,abdeahad2026accurate}. However, these efforts have primarily focused on identifying ambiguities from the perspective of human interpretation. In contrast, our preliminary findings suggest that not all uncertainties perceived by humans necessarily translate into uncertainty for the model.%In contrast, our work studies requirement uncertainty from a model-centered perspective. Instead of assuming that human-perceived ambiguity necessarily leads to generation failure, we characterize requirement uncertainty according to observable inconsistency in repeated code generation.

Recent work on prompt robustness further shows that LLMs are sensitive to natural-language variations, including paraphrasing, formatting changes, noise, and other perturbations~\cite{morris2020textattack,zhu2024promptrobust,chen2026nlperturbator}. These studies provide evidence that small changes in natural-language inputs may affect model outputs. However, our goal is not only to evaluate robustness under perturbations, but to predict them before the downstream tasks. %Our work is related to this line of research because we also examine how requirement variations influence code generation behavior. However, our goal is not only to evaluate robustness under perturbations. We further use generation inconsistency as a signal to quantify requirement uncertainty, predict it before generation, and selectively trigger clarification.

% \subsection{Uncertainty and Reliability in Code Generation}

% Large language models have shown strong capability in code generation, and benchmarks such as HumanEval and MBPP have been widely used to evaluate functional correctness~\cite{humaneval,mbpp}. Subsequent benchmarks and evaluation frameworks further improve test coverage and reliability assessment for generated code~\cite{liu2023evalplus,jain2025livecodebench}. In addition to pass-rate based evaluation, many studies have explored sampling-based generation strategies, self-consistency, and execution-guided selection to improve code generation reliability. These methods typically rely on generating multiple candidate programs and selecting or repairing outputs based on tests or other feedback.

% Our work addresses a complementary problem. Rather than estimating reliability only after code has been generated, we aim to identify risky requirements before the final generation step. To this end, we use repeated generation only in an offline stage to construct model-centered uncertainty signals, including functional instability and structural divergence. A lightweight classifier is then trained to predict uncertainty directly from requirement text. This design allows uncertainty-aware decision making before generation, reducing the need for repeated generation during deployment.

\subsection{Requirements Clarification for Code Generation}
Compared to the extensive work on code generation, research on requirements clarification is still limited. ClarifyGPT~\cite{mu2024clarifygpt} and SpecFix~\cite{jia2025specfix} are two representative approaches. ClarifyGPT iteratively detects ambiguity by monitoring output stability across generations and then asks clarifying questions, while SpecFix revises requirements based on compiler messages, test failures, or execution feedback. Both studies confirm that clarifying requirements benefits code quality. However, they operate after code generation, by regenerating or revising outputs, rather than improving requirement quality before generation. This highlights the need for pre-generation uncertainty identification and mitigation, as early intervention could avoid costly iterative cycles altogether.
%Another related line of work improves code generation through clarification, feedback, or repair. 
%Interactive approaches ask additional questions when the input specification is incomplete or ambiguous, while repair-based approaches revise generated programs using compiler messages, test failures, or execution feedback~\cite{mu2024clarifygpt,jia2025specfix}. These methods show that additional information and iterative refinement can improve the correctness of generated code. However, many existing approaches apply clarification or repair after an initial generation attempt, or use interaction broadly without explicitly deciding which requirements need clarification.

%Our work differs in two aspects. First, we focus on pre-generation clarification: the system predicts whether a requirement is likely to induce uncertain generation before producing the final code. Second, we use uncertainty prediction as a trigger for selective clarification. This avoids unnecessary interaction for low-uncertainty requirements while focusing clarification effort on requirements that are more likely to benefit from additional information. In this sense, our work connects uncertainty quantification, pre-generation prediction, and clarification into a unified workflow for improving the reliability of LLM-based code generation.

